\begin{textsample}{POS Dim 4 – 2019\_c4 – Score 2.00 – 2019\_c4/enade\_2019\_eng\_contr\_b307920d5e.txt}  \label{ex:f4_pos_115}
CAPÍTULO 4 PERCEPÇÃO DA PROVA Nas análises feitas neste capítulo, constam as percepções dos concluintes da Área de Engenharia de Controle e Automação sobre a prova aplicada no Enade/2019.
Essas percepções foram mensuradas por meio de nove questões que avaliaram desde o grau de dificuldade da prova até o tempo gasto para concluí -la.
As percepções sobre a prova foram relacionadas ao desempenho dos estudantes, à Grande Região de funcionamento do curso, à Categoria Administrativa e ao tipo de Organização Acadêmica da IES.
O questionário de percepção da prova encontra -se ao final do Anexo VII, que traz a reprodução do exame.
O desempenho dos estudantes foi classificado em quatro quartos.
Para tanto, esse desempenho foi ordenado de forma ascendente.
O percentil 25, P25, também conhecido como primeiro quartil, é a nota de desempenho que deixa um quarto ( 25 \% ) dos valores observados abaixo e três quartos acima.
A Figura 1 apresenta uma ilustração desse conceito.
O quarto inferior de desempenho é composto pelas notas abaixo do primeiro quartil.
Já o percentil 75, P75, também conhecido como terceiro quartil, é o valor para o qual há três quartos ( 75 \% ) dos dados abaixo e um quarto acima dele.
O quarto superior de desempenho é composto pelas notas iguais ou acima do terceiro quartil.
O percentil 50, P50, também conhecido como mediana, é o valor que divide as notas em dois conjuntos de igual tamanho.
O segundo quarto inclui valores entre o primeiro quartil ( P25 ) e a mediana.
O terceiro quarto contém os valores entre a mediana ( P50 ) e o terceiro quartil ( P75 ).
Vale ressaltar que percentis, quartis e medianas são pontos que, não obrigatoriamente, pertencem ao conjunto original de dados, ao passo que os quartos são subconjuntos dos dados originais.
A seguir, serão apresentados gráficos com resultados selecionados, relativos às nove questões avaliadas por grupos de estudantes.
Nas barras dos gráficos, apresenta -se o percentual de alunos que assinalaram uma das opções ou a soma das porcentagens daqueles que assinalaram duas ( ou três ) delas.
Por exemplo, para as questões 1 e 2, nos gráficos, é apresentada a porcentagem total de participantes que assinalaram as opções ( D ) Difícil e ( E ) Muito difícil.
Em cada barra, foram assinalados, também, os extremos do intervalo de confiança de 95 \% como linhas verticais unidas por uma linha horizontal na forma da letra H maiúscula.
O estimador de um parâmetro com um certo nível de confiança ( e.g. 95,0 \% ) deve conter o parâmetro no intervalo de confiança em 95 \% das vezes.
Na comparação entre os estimadores dos parâmetros de duas classes de uma dada categoria ( e.g.
Norte e Nordeste nas Grandes Regiões, de primeiro e último quarto dentro de desempenho, de IES Públicas e Privadas, ou de Universidades e Faculdades ) associados aos seus respectivos intervalos de confiança, diz -se que não há diferença estatisticamente significativa entre os parâmetros das duas categorias quando há interseção entre os intervalos de confiança e que há diferença, se os intervalos de confiança são disjuntos ( para maiores informações vide Glossário ).
Nas Tabelas no Anexo II, são apresentados os valores absolutos e a distribuição percentual18 das alternativas válidas das nove questões, segundo o mesmo recorte de desempenho dos alunos e Grande Região de funcionamento do curso, de Categoria Administrativa e do tipo de Organização Acadêmica da IES. 4.1 GRAU DE DIFICULDADE DA PROVA 4.1.1 Componente de Formação Geral Ao avaliarem Qual o grau de dificuldade desta prova na parte de Formação Geral? ( Questão 1 ), 28,9 \% do grupo de inscritos e presentes optaram pelas alternativas Difícil ou Muito difícil.
Entretanto, para mais da metade dos estudantes ( 58,1 \% ), o Componente de Formação Geral da prova foi considerado com grau de dificuldade Médio ( Gráfico 4.1, Gráfico 4.2 e, no Anexo II, a Tabela II.1 ).
O percentual de estudantes que consideraram a prova como Difícil ou Muito difícil foi maior na região Sudeste, onde a proporção foi de 29,6 \%, enquanto a de menor incidência foi a Centro-Oeste, com 26,3 \%.
No Gráfico 4.1, é possível observar que as diferenças entre as regiões não são estatisticamente significativas.
Nas Grandes Regiões, a proporção de presentes à prova que consideraram o Componente de Formação Geral como sendo de grau de dificuldade Médio esteve entre 56,1 \%, na região Sul e 64,4 \%, na região Norte.
O percentual de alunos que consideraram o componente de Formação Geral da prova Difícil ou Muito difícil foi decrescente em função dos quartos de desempenho : 39,3 \%, no primeiro quarto e 19,5 \%, no quarto, grupo de melhor desempenho na prova.
Nos quartos de desempenho intermediários, a proporção de alunos que consideraram a prova Difícil ou Muito difícil foi de 31,0 \%, no segundo quarto e de 25,9 \%, no terceiro quarto.
As diferenças entre todos os quartos de desempenho são estatisticamente significativas, com exceção da diferença entre o segundo e o terceiro quartos.
Para todos os quartos de desempenho, a alternativa modal para essa pergunta foi a Médio, com 51,1 \% e 61,5 \% dos respondentes nos quartos extremos, primeiro e quarto, respectivamente.
O percentual de estudantes que consideraram o componente de Formação Geral da prova como Difícil ou Muito difícil foi maior para os estudantes de Faculdades, para os quais a proporção foi de 39,7 \%, enquanto os CEFET/IF foram a Organização Acadêmica com a menor incidência, com 22,9 \%, com diferença estatisticamente significativa entre os dois tipos de organização.
No Gráfico 4.3, é possível observar que as diferenças de Universidades ( 24,2 \% ) e CEFET/IF para Centros Universitários ( 36,6 \% ) e Faculdades são estatisticamente significativas.
Nas Organizações Acadêmicas, a proporção de presentes à prova que consideraram o Componente de Formação Geral como sendo de grau de dificuldade Médio esteve entre 53,1 \%, nas Faculdades e 62,4 \% nos CEFET/IF ( ver também a Tabela II.2 no Anexo II ).
O percentual de estudantes que consideraram o componente de Formação Geral da prova como Difícil ou Muito difícil foi maior para os estudantes de cursos de IES Privadas ( 36,3 \% ), uma diferença estatisticamente significativa da proporção em IES Públicas ( 20,3 \% ).
Já a proporção de presentes à prova que consideraram o Componente de Formação Geral como sendo de grau de dificuldade Médio foi 61,6 \% para os de IES Públicas e 55,0 \% para os de IES Privadas ( ver Gráfico 4.4 ). 4.1.2 Componente de Conhecimento Específico Ao responderem à Questão 2 – Qual o grau de dificuldade desta prova na parte de Componente Específico? – 65,5 \% do grupo de estudantes classificaram -na como Difícil ou Muito difícil.
Além disso, o Componente modal de Conhecimento Específico da prova foi considerado com grau de dificuldade Médio por 31,9 \% dos alunos ( Gráfico 4.5, Gráfico 4.6, e, no Anexo II, a Tabela II.3 ).
A análise das respostas dos estudantes quanto ao grau de dificuldade do Componente de Conhecimento Específico da prova, agregado por Grande Região, mostra que a diferença entre a menor e as duas maiores proporções de alunos que a avaliaram como Difícil ou Muito difícil é estatisticamente significativa, sendo a menor, na região Nordeste ( 58,4 \% ), e a maior na região Norte ( 70,6 \% ).
O percentual de alunos que classificaram o grau de dificuldade como Médio, no Componente de Conhecimento Específico, variou de 29,4 \% a 39,7 \%, para as regiões Norte e Nordeste, respectivamente.
Considerando -se a avaliação da dificuldade das questões do Componente de Conhecimento Específico da prova, de acordo com o desempenho dos estudantes, observa -se que há diferença estatisticamente significativa dos resultados entre o primeiro quarto e o terceiro e último quartos de desempenho, assim como entre o último quarto e os demais quartos.
A proporção dos que classificaram a parte específica como Difícil ou Muito difícil variou de 53,8 \% ( último quarto ) a 72,3 \% ( primeiro quarto ).
As demais proporções dessa \textbf{reposta} foram de 69,3 \% e 66,7 \%, nos segundo e terceiro quartos, respectivamente.
Já a proporção dos que responderam que o grau de dificuldade das questões do Componente de Conhecimento Específico da prova foi Médio variou de 25,0 \%, no primeiro quarto a 41,8 \%, no último.
Considerando -se a avaliação da dificuldade das questões do Componente de Conhecimento Específico da prova, o percentual de estudantes que as consideraram como Difícil ou Muito difícil foi maior para os estudantes de Faculdades, para os quais a proporção foi de 73,4 \%, enquanto os CEFET/IF foram a Organização Acadêmica com a menor incidência, com 62,0 \%.
No Gráfico 4.7, é possível observar que as diferenças de Universidades ( 62,5 \% ) e CEFET/IF para Faculdades e Centros Universitários ( 69,1 \% ) são estatisticamente significativas.
Nas Organizações Acadêmicas, a proporção de presentes à prova que consideraram o Componente de Conhecimento Específico como sendo de grau de dificuldade Médio esteve entre 24,9 \%, nas Faculdades e 35,4 \%, nos CEFET/IF ( ver também Tabela II.4 no Anexo II ).
O percentual de estudantes que consideraram o Componente Específico da prova como Difícil ou Muito difícil foi maior para os estudantes de cursos de IES Privadas ( 72,0 \% ), uma diferença estatisticamente significativa da proporção em IES Públicas ( 58,0 \% ).
Já a proporção de presentes à prova que consideraram este componente como sendo de grau de dificuldade Médio foi 38,6 \% para os de IES Públicas e 26,2 \% para os de IES Privadas ( ver Gráfico 4.8 e Tabela II.4 no Anexo II ). 4.2 EXTENSÃO DA PROVA EM RELAÇÃO AO TEMPO TOTAL Indagados quanto à extensão da prova, em relação ao tempo total oferecido para a sua resolução ( Questão 3 ), os estudantes apontaram, com maior incidência, a alternativa que considerava a extensão adequada, para todas as agregações consideradas ( Gráfico 4.9, Gráfico 4.10, e, no Anexo II, a Tabela II.5 ).
O percentual de alunos que responderam ser a extensão da prova Adequada foi de 56,2 \%.
Já 34,0 \% dos inscritos presentes consideraram que a prova foi Longa ou Muito longa, e 9,8 \% a avaliaram como Curta ou Muito curta.
Dentre as Grandes Regiões, a proporção daqueles que avaliaram a prova como Longa ou Muito longa, em relação ao tempo total destinado à sua resolução, variou de 24,9 \%, na região Norte até 40,4 \%, na região Centro-Oeste.
Há diferença estatisticamente significativa entre a região Norte e todas as demais regiões.
Considerando -se o desempenho dos alunos, nota -se, ainda, que 57,2 \% consideraram a extensão da prova Adequada, no quarto de desempenho inferior e 54,1 \%, no de melhor desempenho ( quarto superior ).
Nos quartos intermediários, essa proporção foi de 56,7 \%, no segundo e no terceiro quartos, proporções decrescentes com o desempenho.
No Gráfico 4.10, pode -se constatar que não há diferença estatisticamente significativa entre as proporções de estudantes que consideraram a prova Longa ou Muito longa em função dos quartos de desempenho, sendo a maior proporção dessa resposta no último quarto de desempenho ( 37,1 \% ), e a menor, no primeiro quarto ( 32,5 \% ), mas não caracterizando uma tendência.
Considerando -se a extensão da prova, o percentual de estudantes que avaliaram a prova como Longa ou Muito longa, em relação ao tempo total destinado à sua resolução foi maior para os estudantes de CEFET/IF, para os quais a proporção foi de 36,4 \%, enquanto as Faculdades foram a Organização Acadêmica com a menor incidência, com 33,2 \%.
No Gráfico 4.11, é possível observar que as diferenças entre CEFET/IF, Faculdades, Centros Universitários ( 33,9 \% ) e Universidades ( 33,6 \% ) não são estatisticamente significativas.
Nas Organizações Acadêmicas, a proporção de presentes à prova que consideraram a sua extensão como adequada, em relação ao tempo total destinado à sua resolução esteve entre 52,5 \%, nos CEFET/IF e 57,6 \%, nas Universidades ( ver também Tabela II.6 no Anexo II ).
O percentual de estudantes que consideraram a prova como Longa ou Muito longa, em relação ao tempo total destinado à sua resolução foi maior para os estudantes de cursos de IES Públicas ( 34,8 \% ), sem diferença estatisticamente significativa da proporção em IES Privadas ( 33,4 \% ).
Já a proporção de presentes à prova que consideraram a extensão como Adequada foi 56,0 \% para os de IES Públicas e 56,4 \% para os de IES Privadas ( ver Gráfico 4.12 e Tabela II.6 no Anexo II ). 4.3 COMPREENSÃO DOS ENUNCIADOS DAS QUESTÕES 4.3.1 Componente de Formação Geral Em relação aos enunciados das questões do Componente de Formação Geral ( Questão 4 ), as opiniões foram positivas, já que 74,1 \% dos alunos avaliados consideraram todos ou a maioria dos enunciados das questões claros e objetivos ( Gráfico 4.13, Gráfico 4.14, e, no Anexo II, a Tabela II.7 ).
Na análise regional, a percentagem de estudantes que avaliaram que todos ou a maioria dos enunciados das questões do Componente de Formação Geral estavam claros e objetivos variou de 69,3 \%, na região Sul a 75,9 \%, na região Sudeste.
A análise das percepções dos estudantes sobre a clareza e objetividade dos enunciados permite afirmar que todos, ou a maioria dos enunciados de questões relativas ao Componente de Formação Geral, foram considerados claros e objetivos para a maior parte dos respondentes ( maior ou igual a 69,3 \%, em todas as regiões e maior ou igual a 72,2 \% para todos os quartos de desempenho ).
Segundo o desempenho, observa -se que a proporção dos que emitiram essa opinião tende a crescer com o desempenho, sem diferenças estatisticamente significativas entre os quartos de desempenho.
No quarto superior, a clareza e a objetividade de todos ou da maioria dos enunciados das questões foram percebidas por 75,3 \%.
A proporção dos presentes que consideraram que todos ou a maioria dos enunciados das questões do Componente de Formação Geral estavam claros e objetivos foi maior para os estudantes de Faculdades, para os quais a proporção foi de 76,3 \%, enquanto os CEFET/IF foram a Organização Acadêmica com a menor incidência, com 70,9 \%.
No Gráfico 4.15, é possível observar que as diferenças entre CEFET/IF, Faculdades, Centros Universitários ( 74,1 \% ) e Universidades ( 74,5 \% ) não são estatisticamente significativas ( ver também Tabela II.8 no Anexo II ).
O percentual de estudantes que consideraram que todos ou a maioria dos enunciados das questões do Componente de Formação Geral estavam claros e objetivos foi maior para os estudantes de cursos de IES Privadas ( 75,2 \% ), sem diferença estatisticamente significativa da proporção em IES Públicas ( 73,0 \% ).
Um maior detalhamento pode ser encontrado no Gráfico 4.16 e na Tabela II.8 no Anexo II. 4.3.2 Componente de Conhecimento Específico Em relação aos enunciados das questões do Componente de Conhecimento Específico da prova, para 80,8 \% dos estudantes avaliados da Área de Engenharia de Controle e Automação, a clareza e a objetividade ( Questão 5 ) estavam presentes em todas ou na maioria das questões ( Gráfico 4.17 a Gráfico 4.20, e no Anexo II, as Tabelas II.9 e II.10 ).
A maioria dos estudantes de todas as Grandes Regiões brasileiras considerou claros e objetivos todos ou a maioria dos enunciados das questões do Componente de Conhecimento Específico da prova, percentual sempre maior ou igual a 78,8 \%.
A diferença entre as regiões não é estatisticamente significativa.
A proporção de estudantes que consideraram os enunciados das questões claros e objetivos apresenta uma tendência crescente em relação ao aumento de desempenho : mais elevada no quarto superior ( 86,7 \% ), se comparada ao quarto inferior de desempenho ( 72,9 \% ).
As diferenças entre os quartos de desempenho são estatisticamente significativas, com exceção da diferença entre o terceiro e o quarto quartos.
A proporção dos presentes que consideraram que todos ou a maioria dos enunciados das questões do Componente de Conhecimento Específico estavam claros e objetivos foi maior para os estudantes de CEFET/IF, para os quais a proporção foi de 84,5 \%, enquanto as Faculdades foram a Organização Acadêmica com a menor incidência, com 76,2 \%.
No Gráfico 4.19, é possível observar que as diferenças entre Faculdades e Universidades ( 82,3 \% ), bem como entre os CEFET/IF e as Faculdades e Centros Universitários são estatisticamente significativas ( ver também Tabela II.10 no Anexo II ).
O percentual de estudantes que consideraram que todos ou a maioria dos enunciados das questões do Componente de Conhecimento Específico estavam claros e objetivos foi maior para os estudantes de cursos de IES Públicas ( 84,1 \% ), uma diferença estatisticamente significativa da proporção em IES Privadas ( 78,0 \% ).
Ver Gráfico 4.20 e Tabela II.10 no Anexo II para um maior detalhamento. 4.4 SUFICIÊNCIA DAS INFORMAÇÕES/INSTRUÇÕES FORNECIDAS Ao avaliarem as informações/instruções fornecidas para a resolução das questões ( Questão 6 ), 83,6 \% dos respondentes da Área de Engenharia de Controle e Automação de todo o Brasil afirmaram que estas eram até excessivas ou suficientes em todas ou na maioria das questões ( Gráfico 4.21, Gráfico 4.22, e, no Anexo II, a Tabela II.11 ).
Quanto à distribuição de respondentes pelas Grandes Regiões, observa -se que a proporção de estudantes que consideraram as informações/instruções fornecidas até excessivas ou suficientes em todas ou na maioria das questões foi sempre superior ou igual a 74,6 \%, chegando a 86,0 \%, na região Sul.
Existem diferenças estatisticamente significativas entre a região Norte, a menor proporção, e as regiões Sudeste e Sul, as duas maiores proporções.
Levando -se em conta o desempenho dos participantes, observa -se uma tendência crescente das proporções de participantes que avaliaram as informações/instruções como até excessivas ou suficientes em todas ou na maioria das questões, havendo diferenças estatisticamente significativas entre todos os quartos de desempenho, com exceção daquela entre o segundo e o terceiro quartos, como mostrado no Gráfico 4.22.
O percentual foi mais elevado no quarto superior ( 90,9 \% ), percentual superior à média nacional ( 83,6 \% ).
Já no quarto inferior, a suficiência das informações/instruções declarada como até excessiva, em todas ou na maioria das questões foi percebida por 75,7 \% dos respondentes.
Considerando -se as informações/instruções fornecidas para a resolução das questões ( Questão 6 ), a proporção que considerou que estas eram até excessivas ou suficientes em todas ou na maioria das questões foi maior para os estudantes de CEFET/IF, para os quais a proporção foi de 87,1 \%, enquanto as Faculdades foram a Organização Acadêmica com a menor incidência, com 78,3 \%.
No Gráfico 4.23, é possível observar que as diferenças entre as duas maiores proporções, em Universidades ( 85,9 \% ) e CEFET/IF ( 87,1 \% ), e as duas menores, em Centros Universitários ( 79,3 \% ) e Faculdades ( 78,3 \% ) são estatisticamente significativas ( ver também Tabela II.12 no Anexo II ).
O percentual de estudantes que consideraram que as informações/instruções fornecidas para a resolução das questões eram até excessivas ou suficientes em todas ou na maioria das questões foi maior para os estudantes de cursos de IES Públicas ( 88,3 \% ), uma diferença estatisticamente significativa da proporção em IES Privadas ( 79,5 \% ) ( ver Gráfico 4.24 e Tabela II.12 no Anexo II ). 4.5 DIFICULDADE ENCONTRADA AO RESPONDER À PROVA Indagados sobre as dificuldades com as quais se depararam ao responder à prova ( Questão 7 ), 22,6 \% dos estudantes apontaram o Desconhecimento do conteúdo.
Para 46,5 \%, a Forma diferente de abordagem do conteúdo foi indicada como dificuldade.
Já a Falta de motivação para fazer a prova foi a dificuldade apontada por 17,8 \% dos respondentes.
Considerando -se todo o Brasil, 8,5 \% dos respondentes afirmaram que não tiveram qualquer tipo de dificuldade para responder à prova ( Tabelas II.13 e II.14 no Anexo II ).
Nos Gráficos 4.25 a 4.28, são apresentados os percentuais de estudantes que apontaram o Desconhecimento do conteúdo como dificuldade ao responder à prova.
Na análise por Grandes Regiões, o percentual de inscritos e presentes que apontaram o Desconhecimento do conteúdo como dificuldade ao responder à prova não superou 28,8 \%.
Os percentuais variaram de 21,6 \%, na região Sudeste a 28,8 \%, na região Norte.
A Forma diferente de abordagem do conteúdo foi a escolha modal dos estudantes, com percentuais que variaram de 40,8 \% ( região Sul ) a 48,4 \% ( região Sudeste ).
O percentual de alunos que citaram a Falta de motivação como dificuldade variou de 15,8 \% ( região Norte ) a 19,5 \% ( região Sul ).
Os que declararam não ter qualquer tipo de dificuldade para responder à prova variaram de 5,1 \%, na região Norte a 9,4 \%, nas regiões Nordeste e Sul.
Em relação aos quartos de desempenho, o Desconhecimento do conteúdo foi a opção escolhida por 21,6 \% dos estudantes do quarto inferior e por 24,8 \% do quarto superior, caracterizando uma tendência crescente, com exceção do primeiro para o segundo quarto.
Não são observadas diferenças estatisticamente significativas entre os quartos de desempenho.
A alternativa modal para os alunos, quando agregados pelos quartos de desempenho, para a dificuldade encontrada, foi causada pela Forma diferente de abordagem do conteúdo : 50,4 \%, no quarto inferior e 42,5 \%, do quarto superior escolheram essa alternativa.
Na análise por Organização Acadêmica, o percentual de inscritos e presentes que apontaram o Desconhecimento do conteúdo como dificuldade ao responder à prova não superou 24,6 \%.
Os percentuais variaram de 19,9 \%, para CEFET/IF a 24,6 \%, nas Universidades.
A Forma diferente de abordagem do conteúdo foi a escolha modal dos estudantes, com percentuais que variaram de 40,0 \% ( Universidades ) a 59,1 \% ( Faculdades ).
O percentual de alunos que citaram a Falta de motivação como dificuldade variou de 8,7 \% ( Centros Universitários ) a 24,9 \% ( CEFET/IF ).
Os que declararam não ter qualquer tipo de dificuldade para responder à prova variaram de 6,8 \%, nos CEFET/IF a 9,1 \%, nos Centros Universitários.
No Gráfico 4.27, é possível observar que as diferenças entre os tipos de Organização Acadêmica não são estatisticamente significativas ( ver também Tabela II.14 no Anexo II ).
Em relação às Categorias Administrativas, o Desconhecimento do conteúdo foi a opção escolhida por 23,3 \% dos estudantes nas IES Públicas e por 22,0 \% nas IES Privadas.
Não são observadas diferenças estatisticamente significativas entre estas duas categorias.
A alternativa modal para os alunos, quando agregados pelas Categorias Administrativas, para a dificuldade encontrada, foi causada pela Forma diferente de abordagem do conteúdo : 37,0 \%, para as IES Públicas e 54,7 \%, para as IIES Privadas ( ver Gráfico 4.28 e Tabela II.14 no Anexo II ). 4.6 CONTEÚDOS DAS QUESTÕES OBJETIVAS DA PROVA Ao analisar os conteúdos das questões objetivas da prova ( Questão 8 ), um percentual muito pequeno dos estudantes avaliados, apenas 4,1 \%, afirmaram que não estudou ainda a maioria desses conteúdos ( Gráficos 4.29 a Gráfico 4.32, e as Tabelas II.15 e II.16 no Anexo II ).
A maioria ( 60,0 \% ) afirmou ter estudado e aprendido muitos ou todos os conteúdos.
Na análise por Grande Região, a proporção de respondentes que escolheram a opção não estudou ainda a maioria desses conteúdos foi pequena.
Observa -se que, nas regiões Sudeste ( 4,0 \% ), Sul ( 2,9 \% ) e Centro-Oeste ( 3,3 \% ), as proporções foram menores do que a média nacional ( 4,1 \% ).
Nas regiões Norte ( 6,2 \% ) e Nordeste ( 5,7 \% ), a situação foi inversa : uma proporção mais elevada.
Não são observadas diferenças estatisticamente significativas entre as regiões.
Em todas as regiões, a maior parte dos presentes afirmou ter estudado e aprendido muitos ou todos os conteúdos, com proporções variando entre 41,2 \%, na região Norte a 65,4 \%, na região Sul.
Considerando -se separadamente as opiniões de estudantes dos quatro quartos de desempenho, observa -se que, no quarto inferior, 6,7 \% ofereceram como resposta que não estudou ainda a maioria desses conteúdos, sendo 1,6 \% os do quarto superior com a mesma resposta.
A diferença dentre os alunos que optaram por esse motivo de dificuldade nos quartos inferiores, em relação aos quartos superiores de desempenho, é estatisticamente significativa.
Tendo -se em conta o quarto superior, 74,3 \% dos alunos afirmaram ter estudado e aprendido muitos ou todos os conteúdos.
No outro extremo, no primeiro quarto, 47,9 \% optaram pelas mesmas categorias.
As proporções são crescentes com o desempenho.
Na análise por Organização Acadêmica, a proporção de respondentes que escolheram a opção não estudou ainda a maioria desses conteúdos foi pequena.
Observa -se que, nas Faculdades ( 5,8 \% ) e nos Centros Universitários ( 4,6 \% ), as proporções foram maiores do que a média nacional ( 4,1 \% ).
Nos CEFET/IF ( 1,8 \% ) e nas Universidades ( 4,0 \% ) a situação foi inversa : uma proporção mais baixa.
No Gráfico 4.31, é possível observar que as diferenças entre CEFET/IF e as demais organizações acadêmicas são estatisticamente significativas ( ver também Tabela II.16 no Anexo II ).
Em todos os tipos de Organização Acadêmica, a maioria dos presentes afirmou ter estudado e aprendido muitos ou todos os conteúdos, com proporções variando entre 52,0 \%, nas Faculdades a 66,7 \%, nos CEFET/IF.
Em relação às Categorias Administrativas, observa -se que, nas IES Públicas, 2,8 \% ofereceram como resposta que não estudou ainda a maioria desses conteúdos, sendo 5,2 \% nas IES Privadas com a mesma resposta.
A diferença entre os alunos que optaram por esse motivo de dificuldade nas duas categorias é estatisticamente significativa.
Considerando -se as IES Públicas, 66,0 \% dos alunos afirmaram ter estudado e aprendido muitos ou todos os conteúdos.
Na outra categoria, 54,9 \% optaram pelas mesmas categorias ( ver Gráfico 4.32 e Tabela II.16 no Anexo II ). 4.7 TEMPO GASTO PARA CONCLUIR A PROVA Ao responderem sobre o tempo de conclusão da prova ( Questão 9 ), mais de três quartos dos estudantes ( 75,6 \% ) afirmaram ter gasto Entre duas e quatro horas ( Gráfico 4.33 a Gráfico 4.36 e, no Anexo II, as Tabelas II.17 e II.18 ).
Considerando -se as cinco Grandes Regiões brasileiras, os percentuais dos que utilizaram Entre duas e quatro horas para finalizar a prova nas regiões Nordeste ( 71,3 \% ), Sul ( 73,9 \% ) e Centro-Oeste ( 73,7 \% ) foram inferiores ao percentual nacional.
Nas demais Grandes Regiões, o percentual de alunos que dispensaram Entre duas e quatro horas para concluir a prova foi maior : 79,4 \% na região Norte e 76,6 \%, na região Sudeste, como mostrado no Gráfico 4.33.
A diferença entre as regiões não é estatisticamente significativa.
Analisadas as alternativas escolhidas pelos estudantes que se situam nos diferentes quartos de desempenho, observa -se que uma proporção maior de participantes nos quartos superiores declarou ter gasto Entre duas e quatro horas para concluir a prova quando comparada com a proporção de participantes nos quartos inferiores, uma situação que apenas não ocorre do terceiro para o último quarto de desempenho.
Há diferença estatisticamente significativa entre o primeiro e o terceiro quarto de desempenho, respectivamente, a menor e a maior proporção de participantes que escolheram essa alternativa.
Na análise por Organização Acadêmica, os percentuais dos que utilizaram Entre duas e quatro horas para finalizar a prova, apenas, nas Universidades ( 72,9 \% ) foram inferiores ao percentual nacional.
Nos demais tipo de Organização Acadêmica, o percentual de alunos que dispensaram Entre duas e quatro horas para concluir a prova foi mais elevado que a média : 76,2 \%, nos CEFET/IF, 79,4 \%, nas Faculdades e 78,6 \%, nos Centros Universitários, como mostrado no Gráfico 4.35.
As diferenças de Universidades em relação às Faculdades e Centros Universitários são estatisticamente significativas.
Em relação às Categorias Administrativas, observa -se que, nas IES Públicas, 72,3 \% declararam ter gasto Entre duas e quatro horas para concluir a prova, sendo 78,3 \% nas IES Privadas com a mesma resposta.
A diferença entre os alunos que optaram por esse motivo de dificuldade nas duas categorias é estatisticamente significativa ( ver Gráfico 4.36 e Tabela II.18 no Anexo II ).

% matched lemmas: repor
\end{textsample}
